\subsection{Stochastic and amortized variational inference}
\label{subsec:amortized-vi}

\subsubsection{From mean-field factors to shared encoders}

Mean-field VI introduces local variational objects for each observation.
If $x_1,\dots,x_n$ are observations with corresponding latent variables $z_1,\dots,z_n$, a classical approximation might introduce one local parameter $\lambda_i$ for each $x_i$.
That can be effective, but it scales poorly:
every new observation comes with a fresh inference subproblem.
In the Bayesian factor-analysis example from \Cref{ex:mfvi-bfa}, this means storing one Gaussian factor \(q_i(z_i)\) for every datum together with a global factor for the loading matrix \(W\).

Amortized variational inference replaces those per-datum free parameters by a shared map.
Instead of optimizing \(\lambda_i\) separately for every \(x_i\), we learn parameters \(\phi\) and set
\[
  q_\phi(z_i\mid x_i)
\]
directly as a function of the observation.
In factor analysis, a natural amortized choice is a Gaussian encoder
\[
  q_\phi(z_i\mid x_i)
  =
  \mathcal{N}\!\bigl(\mu_\phi(x_i),\Sigma_\phi(x_i)\bigr),
\]
which tries to predict the posterior summaries that CAVI would otherwise compute separately for each \(i\).
If the model still contains a global latent block such as \(W\), we keep a separate factor \(q_\xi(W)\) and amortize only the local factors:
\[
  q_{\xi,\phi}(W,z_{1:n})
  =
  q_\xi(W)\prod_{i=1}^n q_\phi(z_i\mid x_i).
\]
This is the meaning of \emph{amortization}: we pay the cost of learning an inference rule during training and then reuse that rule across data points.
When there is no separate global latent block, and the generative model has parameters \(\theta\) that factorize across observations as
\[
  p_\theta(x_{1:n},z_{1:n})
  =
  \prod_{i=1}^n p_\theta(x_i,z_i),
\]
then the training objective is the summed ELBO
\begin{equation}
  \mathcal{J}(\theta,\phi)
  \coloneqq
  \sum_{i=1}^n \mathcal{L}_i(\theta,\phi),
  \qquad
  \mathcal{L}_i(\theta,\phi)
  \coloneqq
  \E_{q_\phi(z_i\mid x_i)}
  \bigl[\log p_\theta(x_i,z_i) - \log q_\phi(z_i\mid x_i)\bigr].
  \label{eq:amortized-elbo}
\end{equation}
Stochastic-gradient VI is the optimization strategy that makes this practical: we estimate gradients of \eqref{eq:amortized-elbo} with Monte Carlo samples and minibatches rather than with closed-form coordinate updates.
\Cref{fig:vi-amortization} illustrates the difference.
This encoder--decoder organization is the modern VAE setup \citep[\S2.3]{kingma_welling_2014_aevb}.

\begin{figure}[H]
  \centering
  \begin{tikzpicture}[>=Stealth,font=\footnotesize]
  \path[use as bounding box] (-0.15,-0.1) rectangle (12.95,4.65);

  \begin{scope}
    \node[anchor=west] at (0.25,4.0) {classical local VI};
    \draw[stat221Gray!35,line width=0.4pt] (6.15,0.35) -- (6.15,4.10);

    \foreach \idx/\y/\xl/\ll in {
      1/3.15/$x_1$/$\lambda_1$,
      2/2.20/$x_2$/$\lambda_2$,
      3/1.25/$x_3$/$\lambda_3$
    }{
      \node[draw=stat221Blue!55!black,fill=stat221BlueLight,rounded corners=2pt,minimum width=1.05cm,minimum height=0.62cm] (x\idx) at (1.35,\y) {\xl};
      \node[draw=stat221Gray!60!black,fill=white,rounded corners=2pt,minimum width=1.35cm,minimum height=0.62cm] (l\idx) at (4.65,\y) {\ll};
      \draw[->,stat221Gray!75,line width=0.7pt] (x\idx.east) -- (l\idx.west);
      \node[stat221Gray,anchor=south] at (3.00,\y+0.12) {solve};
    }

    \node[stat221Gray,align=center] at (3.0,0.30) {one separate optimization\\for each observation};
  \end{scope}

  \begin{scope}[shift={(6.65,0)}]
    \node[anchor=west] at (0.35,4.0) {amortized VI};

    \node[draw=stat221Purple!65!black,fill=stat221PurpleLight,rounded corners=3pt,minimum width=2.15cm,minimum height=2.8cm,align=center] (enc) at (3.25,2.2) {shared encoder\\$q_\phi(z\mid x)$};

    \foreach \idx/\y/\xl/\ql in {
      1/3.15/$x_1$/$q_\phi(z\mid x_1)$,
      2/2.20/$x_2$/$q_\phi(z\mid x_2)$,
      3/1.25/$x_3$/$q_\phi(z\mid x_3)$
    }{
      \node[draw=stat221Blue!55!black,fill=stat221BlueLight,rounded corners=2pt,minimum width=1.05cm,minimum height=0.62cm] (rx\idx) at (0.9,\y) {\xl};
      \node[draw=stat221Green!60!black,fill=stat221GreenLight,rounded corners=2pt,minimum width=2.15cm,minimum height=0.62cm] (rq\idx) at (5.85,\y) {\ql};
      \draw[->,stat221Gray!75,line width=0.7pt] (rx\idx.east) -- (enc.west |- 0,\y);
      \draw[->,stat221Gray!75,line width=0.7pt] (enc.east |- 0,\y) -- (rq\idx.west);
    }

    \node[stat221Gray,align=center] at (3.25,0.30) {one learned map reused\\for every observation};
  \end{scope}
\end{tikzpicture}

  \caption{Classical mean-field VI uses one set of local variational parameters per observation.
  Amortized VI replaces those free parameters by a shared encoder $q_\phi(z\mid x)$, so inference for a new datum becomes a forward pass through the learned map.}
  \label{fig:vi-amortization}
\end{figure}

Dividing \eqref{eq:amortized-elbo} by \(n\) gives the same optimizer, and the corresponding population form replaces the empirical sum by
\[
  \E_{p_{\mathrm{data}}(x)}\!\bigl[\mathcal{L}(q_\phi(\cdot\mid x);x)\bigr].
\]
We keep the finite-sample form here because the main issue in this chapter is stochastic optimization over a given data set.

For the Bayesian factor-analysis model, the amortized family \(q_{\xi,\phi}(W,z_{1:n})\) leads to
\[
  \mathcal{J}(\theta,\xi,\phi)
  =
  \E_{q_{\xi,\phi}}
  \left[
    \log p_\theta(W) - \log q_\xi(W)
    +
    \sum_{i=1}^n
    \bigl(
      \log p_\theta(z_i)
      +
      \log p_\theta(x_i\mid z_i,W)
      -
      \log q_\phi(z_i\mid x_i)
    \bigr)
  \right].
\]
With Gaussian priors and Gaussian observation noise, this simplifies to
\begin{equation}
  \mathcal{J}(\theta,\xi,\phi)
  =
  C
  -
  \operatorname{KL}\!\bigl(q_\xi(W)\,\|\,p_\theta(W)\bigr)
  -
  \sum_{i=1}^n
  \operatorname{KL}\!\bigl(q_\phi(z_i\mid x_i)\,\|\,p_\theta(z_i)\bigr)
  -
  \frac{1}{2\sigma^2}
  \sum_{i=1}^n
  \E_{q_{\xi,\phi}}\norm{x_i-Wz_i}^2,
  \label{eq:amortized-factor-analysis-elbo}
\end{equation}
where \(C\) does not depend on \(\xi\) or \(\phi\).
The Gaussian encoder keeps the local posterior approximation in the same family as the mean-field example, but now \(\mu_\phi(x)\) and \(\Sigma_\phi(x)\) are shared functions rather than \(n\) separate free parameters.
Up to the same additive constant, a minibatch \(B\subseteq\{1,\dots,n\}\) gives the stochastic objective
\begin{equation}
  \hat{\mathcal{J}}_B(\theta,\xi,\phi)
  =
  -\operatorname{KL}\!\bigl(q_\xi(W)\,\|\,p_\theta(W)\bigr)
  +
  \frac{n}{|B|}
  \sum_{i\in B} \ell_i(\theta,\xi,\phi),
  \label{eq:amortized-factor-analysis-minibatch}
\end{equation}
where
\[
  \ell_i(\theta,\xi,\phi)
  \coloneqq
  \E_{q_\xi(W)q_\phi(z_i\mid x_i)}\!\bigl[\log p_\theta(x_i\mid z_i,W)\bigr]
  -
  \operatorname{KL}\!\bigl(q_\phi(z_i\mid x_i)\,\|\,p_\theta(z_i)\bigr).
\]
So the global term appears once, while the local reconstruction and KL terms are the ones scaled by \(n/|B|\).
This is the minibatch consequence that makes the factor-analysis example compatible with stochastic-gradient training.

This same example also explains the move from linear latent-variable models to variational autoencoders.
If the loading matrix is treated as an ordinary model parameter rather than a latent variable, the decoder is simply the linear map \(z\mapsto Wz\).
A variational autoencoder keeps the Gaussian latent prior and Gaussian encoder but replaces that linear mean map by a nonlinear decoder \(c_\theta:\R^k\to\R^d\), typically a neural network, and learns \(\theta\) jointly with \(\phi\):
\[
  p_\theta(x\mid z)
  =
  \mathcal{N}\!\bigl(c_\theta(z),\sigma^2 I_d\bigr),
  \qquad
  q_\phi(z\mid x)
  =
  \mathcal{N}\!\bigl(\mu_\phi(x),\Sigma_\phi(x)\bigr).
\]
For one datum the ELBO then becomes
\[
  \mathcal{L}_i(\theta,\phi)
  =
  C
  -
  \frac{1}{2\sigma^2}
  \E_{q_\phi(z_i\mid x_i)}\norm{x_i-c_\theta(z_i)}^2
  -
  \operatorname{KL}\!\bigl(q_\phi(z_i\mid x_i)\,\|\,p_\theta(z_i)\bigr).
\]
So the conceptual move from factor analysis to a VAE is small:
the latent variable \(z\) stays in place, the encoder stays Gaussian, and only the linear decoder \(Wz\) is replaced by a flexible nonlinear map \citep[\S2.3]{kingma_welling_2014_aevb}.
What changes analytically is that the linear-Gaussian closed forms disappear twice.
First, the reconstruction expectation is no longer available in closed form once \(Wz\) is replaced by \(c_\theta(z)\).
Second, writing the ELBO gradient formally still leaves an intractable expectation under \(q_\phi(z_i\mid x_i)\).
That is why the pathwise reparameterization step below is not a cosmetic trick.
It is the mechanism that turns this nonlinear latent-variable model back into a trainable stochastic optimization problem.

\subsubsection{Reparameterization}
\label{subsubsec:reparameterization-trick}

For continuous latent variables, the Gaussian encoder from the factor-analysis example already has the right form for the reparameterization trick.
The cleanest route is to rewrite the sampled latent variable as a differentiable transformation of parameter-free noise.
It converts the gradient of an expectation under \(q_\phi\) into an ordinary pathwise derivative through a fixed source of noise, which is the main computational idea behind stochastic variational autoencoders \citep[\S2.3]{kingma_welling_2014_aevb} and stochastic backpropagation \citep[\S3]{rezende_mohamed_wierstra_2014_stochastic_backpropagation}.
Suppose we can sample from the variational family by drawing noise $\varepsilon$ from a fixed base distribution and applying a differentiable transformation,
\begin{equation}
  z = T_\phi(x,\varepsilon),
  \qquad
  \varepsilon \sim s(\varepsilon),
  \label{eq:reparameterization-map}
\end{equation}
where the law of $T_\phi(x,\varepsilon)$ is exactly $q_\phi(z\mid x)$.
Then the ELBO term for one observation can be rewritten as
\begin{equation}
  \mathcal{L}_i(\theta,\phi)
  =
  \E_{\varepsilon\sim s}
  \Bigl[
    \log p_\theta\bigl(x_i,T_\phi(x_i,\varepsilon)\bigr)
    -
    \log q_\phi\bigl(T_\phi(x_i,\varepsilon)\mid x_i\bigr)
  \Bigr].
  \label{eq:reparameterized-elbo}
\end{equation}
Now the randomness is independent of the parameters, so gradients with respect to $\theta$ and $\phi$ can be pushed through the sampled latent variable by ordinary backpropagation.
The point is structural: once the latent sample is written as a deterministic function of \((x,\varepsilon)\), the ELBO becomes an expectation over a parameter-free noise source.

The standard example is a Gaussian encoder:
\begin{equation}
  q_\phi(z\mid x)
  =
  \mathcal{N}\!\bigl(\mu_\phi(x),\Sigma_\phi(x)\bigr),
  \qquad
  \Sigma_\phi(x)
  =
  L_\phi(x)L_\phi(x)^\top,
  \qquad
  z
  =
  \mu_\phi(x) + L_\phi(x)\varepsilon,
  \qquad
  \varepsilon\sim\mathcal{N}(0,I).
  \label{eq:gaussian-reparameterization}
\end{equation}
The encoder outputs the mean and a covariance factor, and one latent draw is obtained by transforming standard Gaussian noise.
A diagonal covariance is the common practical specialization in VAEs, but the same pathwise idea applies more generally.
Other common reparameterizable constructions fit the same template: any location-scale family, inverse-CDF transforms when the quantile map is tractable, and richer autoregressive or log-normal style transformations built from differentiable maps of base noise.
This is the basic mechanism behind variational autoencoders.
The encoder produces an approximate posterior, the decoder specifies $p_\theta(x\mid z)$, and the reparameterized ELBO lets us train both with ordinary stochastic gradients.

\begin{algorithm}[H]
  \caption{Stochastic amortized VI with reparameterization}
  \label{alg:amortized-vi}
  \begin{algorithmic}[1]
    \Require Data $x_{1:n}$, generative model $p_\theta(x,z)$, encoder $q_\phi(z\mid x)$ with reparameterization map $T_\phi$, step sizes $\eta_t$
    \For{$t=0,1,2,\dots$}
      \State Sample a minibatch $B\subseteq\{1,\dots,n\}$
      \For{each $i\in B$}
        \State Draw $\varepsilon_i \sim s(\varepsilon)$ and set $z_i \gets T_\phi(x_i,\varepsilon_i)$
      \EndFor
      \State Form the stochastic ELBO estimate
      \[
        \hat{\mathcal{J}}_B(\theta,\phi)
        =
        \frac{n}{|B|}
        \sum_{i\in B}
        \Bigl[
          \log p_\theta(x_i,z_i)
          -
          \log q_\phi(z_i\mid x_i)
        \Bigr]
      \]
      \State Update $(\theta,\phi)$ by a stochastic-gradient step using $\nabla_{\theta,\phi}\hat{\mathcal{J}}_B$
    \EndFor
  \end{algorithmic}
\end{algorithm}

\subsubsection{Minibatch ELBO optimization}
\label{subsubsec:amortized-stochastic-gradients}

The summed objective \eqref{eq:amortized-elbo} is usually optimized with minibatch stochastic gradient ascent.
If $B\subseteq\{1,\dots,n\}$ is a minibatch, the unbiased scaled objective
\begin{equation}
  \hat{\mathcal{J}}_B(\theta,\phi)
  \coloneqq
  \frac{n}{|B|}
  \sum_{i\in B}
  \hat{\mathcal{L}}_i(\theta,\phi)
  \label{eq:minibatch-elbo}
\end{equation}
approximates \eqref{eq:amortized-elbo}, where each term $\hat{\mathcal{L}}_i$ is itself estimated by Monte Carlo.
This is the minibatch VAE training pattern from \citet[\S2.3]{kingma_welling_2014_aevb}.
With $S$ samples $z_i^{(s)}\sim q_\phi(\cdot\mid x_i)$,
\begin{equation}
  \hat{\mathcal{L}}_i(\theta,\phi)
  =
  \frac{1}{S}
  \sum_{s=1}^S
  \Bigl[
    \log p_\theta(x_i,z_i^{(s)})
    -
    \log q_\phi(z_i^{(s)}\mid x_i)
  \Bigr].
  \label{eq:mc-elbo}
\end{equation}
In the factor-analysis objective \eqref{eq:amortized-factor-analysis-elbo}, the same pattern appears with one additional global term:
\[
  -\operatorname{KL}\!\bigl(q_\xi(W)\,\|\,p_\theta(W)\bigr).
\]
This term stays outside the scaled minibatch sum, exactly as in \eqref{eq:amortized-factor-analysis-minibatch}.
In practice these evaluations inherit the numerical issues from \Cref{subsec:stable-evaluation,subsec:geometry-control}:
the ELBO is accumulated on the log scale, transformed latent variables require the Jacobian bookkeeping from \Cref{subsubsec:transformed-densities}, and full-covariance Gaussian families are usually parameterized through the Cholesky constructions from \Cref{subsubsec:cholesky-geometry}.
This is where the variational story reconnects to the optimization part of the course: once the ELBO is estimated on minibatches, its training dynamics are governed by the same stochastic-gradient considerations that motivated momentum and adaptive methods earlier.

\subsubsection{Score-function gradients}

When reparameterization is unavailable, one can still differentiate expectations under \(q_\phi\) by a score-function, or REINFORCE, identity.
The correct formula must account for any explicit \(\phi\)-dependence in the integrand as well.

\begin{proposition}[Score-function differentiation identity]
\label{prop:score-function-gradient}
Let \(q_\phi\) be a differentiable density on a support that does not depend on \(\phi\).
Let \(f_\phi\) be an integrable function such that differentiation can be passed under the integral sign in
\[
  \int f_\phi(z)\,q_\phi(z)\,dz.
\]
Then
\begin{equation}
  \nabla_\phi \E_{q_\phi}[f_\phi(Z)]
  =
  \E_{q_\phi}\!\bigl[\nabla_\phi f_\phi(Z)\bigr]
  +
  \E_{q_\phi}\!\bigl[f_\phi(Z)\,\nabla_\phi \log q_\phi(Z)\bigr].
  \label{eq:score-function-gradient}
\end{equation}
If \(f_\phi=f\) does not depend on \(\phi\), this reduces to
\[
  \nabla_\phi \E_{q_\phi}[f(Z)]
  =
  \E_{q_\phi}\!\bigl[f(Z)\,\nabla_\phi \log q_\phi(Z)\bigr].
\]
\end{proposition}
\begin{proof}
Differentiate under the integral sign:
\[
  \nabla_\phi \E_{q_\phi}[f_\phi(Z)]
  =
  \nabla_\phi \int f_\phi(z)\,q_\phi(z)\,dz
  =
  \int \nabla_\phi f_\phi(z)\,q_\phi(z)\,dz
  +
  \int f_\phi(z)\,\nabla_\phi q_\phi(z)\,dz.
\]
Using $\nabla_\phi q_\phi(z)=q_\phi(z)\nabla_\phi \log q_\phi(z)$ gives
\[
  \nabla_\phi \E_{q_\phi}[f_\phi(Z)]
  =
  \E_{q_\phi}\!\bigl[\nabla_\phi f_\phi(Z)\bigr]
  +
  \int f_\phi(z)\,q_\phi(z)\,\nabla_\phi \log q_\phi(z)\,dz
  =
  \E_{q_\phi}\!\bigl[\nabla_\phi f_\phi(Z)\bigr]
  +
  \E_{q_\phi}\!\bigl[f_\phi(Z)\,\nabla_\phi \log q_\phi(Z)\bigr].
\]
This is \eqref{eq:score-function-gradient}.%
\end{proof}

Applied to one-datum ELBO term
\[
  \mathcal{L}_i(\theta,\phi)
  =
  \E_{q_\phi(z_i\mid x_i)}
  \Bigl[
    \log p_\theta(x_i,z_i)-\log q_\phi(z_i\mid x_i)
  \Bigr],
\]
we take
\[
  f_\phi(z_i)
  =
  \log p_\theta(x_i,z_i)-\log q_\phi(z_i\mid x_i).
\]
Then \(\nabla_\phi f_\phi(z_i)=-\nabla_\phi\log q_\phi(z_i\mid x_i)\), and
\[
  \E_{q_\phi}\!\bigl[\nabla_\phi\log q_\phi(Z\mid x_i)\bigr]
  =
  \nabla_\phi \int q_\phi(z\mid x_i)\,dz
  =
  0.
\]
So \eqref{eq:score-function-gradient} yields the usual score-function ELBO gradient
\[
  \nabla_\phi \mathcal{L}_i(\theta,\phi)
  =
  \E_{q_\phi(z_i\mid x_i)}
  \Bigl[
    \bigl(\log p_\theta(x_i,z_i)-\log q_\phi(z_i\mid x_i)\bigr)
    \nabla_\phi \log q_\phi(z_i\mid x_i)
  \Bigr].
\]
This estimator is very general, but it is often noisy in practice because the scalar weight multiplying the score can vary substantially across samples.
For that reason, the reparameterization path above is the default route whenever the latent variables are continuous, and score-function estimators are best viewed as the fallback that keeps the framework applicable more broadly.

Compared with classical mean-field VI, amortization replaces per-datum coordinate optimization by a single shared inference rule.
That rule is approximate, but it scales to large data sets, supports minibatch training, and makes inference for a new datum a single encoder evaluation.
The gain in scalability comes with two familiar risks.
Even when the variational family is expressive enough for each datum separately, a shared encoder can still miss the best local approximation and leave an amortization gap.
If the decoder is strong enough to explain \(x\) without using \(z\) very much, training can drift toward posterior collapse \citep[\S\S3--4]{he_spokoyny_neubig_bergkirkpatrick_2019_lagging_inference}.
